% =============================================================================
% CHƯƠNG 4: TRIỂN KHAI TRÊN MININET
% =============================================================================
\chapter{Triển khai trên Mininet}

\section{Môi trường cài đặt}

\subsection{Yêu cầu hệ thống}

\begin{table}[H]
    \centering
    \caption{Yêu cầu phần mềm và phần cứng}
    \vspace{0.2cm}
    \begin{tabularx}{\textwidth}{|>{\raggedright\arraybackslash}p{4cm}|>{\raggedright\arraybackslash}X|}
        \hline
        \textbf{Thành phần} & \textbf{Phiên bản / Yêu cầu} \\ \hline
        Hệ điều hành & Ubuntu 22.04 LTS (64-bit) \\ \hline
        RAM & Tối thiểu 4 GB (khuyên nghị 8 GB) \\ \hline
        CPU & 2 cores trở lên \\ \hline
        Disk & 20 GB trống \\ \hline
        Python & Python 3.10+ \\ \hline
        Mininet & 2.3.0+ \\ \hline
        Open vSwitch & 2.17+ \\ \hline
        iperf3 & 3.9+ \\ \hline
        iproute2 & Hỗ trợ MPLS (kernel 4.1+) \\ \hline
    \end{tabularx}
    \label{tab:requirements}
\end{table}

\subsection{Hướng dẫn cài đặt}

\textbf{Bước 1: Cập nhật hệ thống}

\begin{lstlisting}[style=bashstyle, caption={Cập nhật hệ thống Ubuntu}]
sudo apt update && sudo apt upgrade -y
\end{lstlisting}

\textbf{Bước 2: Cài đặt các gói cần thiết}

\begin{lstlisting}[style=bashstyle, caption={Cài đặt Mininet và các công cụ hỗ trợ}]
sudo apt install -y \
    mininet \
    openvswitch-switch openvswitch-common \
    python3 python3-pip \
    iperf3 iperf \
    tcpdump net-tools iproute2 \
    iputils-ping traceroute \
    python3-matplotlib python3-numpy \
    curl wget git
\end{lstlisting}

\textbf{Bước 3: Kích hoạt MPLS kernel module}

\begin{lstlisting}[style=bashstyle, caption={Load module MPLS và bật platform labels}]
# Load module MPLS
sudo modprobe mpls_router
sudo modprobe mpls_iptunnel
sudo modprobe mpls_gso

# Kiem tra module da load:
lsmod | grep mpls

# Bat platform labels
echo 'net.mpls.platform_labels = 100000' | \
    sudo tee -a /etc/sysctl.conf
sudo sysctl -p
\end{lstlisting}

\textbf{Bước 4: Kiểm tra Mininet hoạt động}

\begin{lstlisting}[style=bashstyle, caption={Kiểm tra Mininet với topo đơn giản}]
# Test don gian: topo single,3 voi ping
sudo mn --topo single,3 --test ping

# Ket qua mong doi:
# *** Results: 0% dropped (6/6 received)

# Don dep sau test
sudo mn -c
\end{lstlisting}

\section{Cấu trúc file code}

\begin{Verbatim}[fontsize=\small, frame=single, label={Cây thư mục dự án}]
TKM2/
+-- topology/
|   +-- metro_full.py        <- Orchestrator tong the (file chinh)
|   +-- configure_mpls.py    <- Cau hinh MPLS Push/Swap/Pop
|   +-- topo_branch1_flat.py
|   +-- topo_branch2_3tier.py
|   +-- topo_branch3_leafspine.py
|   +-- topo_backbone_mpls.py
+-- scripts/
|   +-- run_measurements.py  <- Do luong tu dong
|   +-- plot_results.py      <- Ve bieu do
|   +-- parse_ping.py
|   +-- parse_iperf.py
|   +-- aggregate_results.py
+-- config/
|   +-- tests.yaml           <- Cau hinh kich ban do luong
+-- results/
|   +-- raw/                 <- Log thu
|   +-- csv/                 <- Du lieu tong hop
|   +-- charts/              <- Bieu do PNG
+-- source_latex/            <- Bao cao LaTeX
+-- README.md
\end{Verbatim}

\section{Mã nguồn chi tiết}

\subsection{Lớp \texttt{LinuxRouter} và \texttt{MetroFullTopo}}

Lớp \texttt{LinuxRouter} kế thừa từ \texttt{Node} của Mininet, bật \texttt{ip\_forward} để node hoạt động như router Linux thật sự:

\begin{lstlisting}[style=pythonstyle, caption={Lớp LinuxRouter -- node hoạt động như IP router}]
class LinuxRouter(Node):
    """Node bat ip_forward -> hoat dong nhu Linux router."""
    def config(self, **params):
        super().config(**params)
        self.cmd('sysctl -w net.ipv4.ip_forward=1')

    def terminate(self):
        self.cmd('sysctl -w net.ipv4.ip_forward=0')
        super().terminate()
\end{lstlisting}

Lớp \texttt{MetroFullTopo} xây dựng toàn bộ topology theo phương pháp modular:

\begin{lstlisting}[style=pythonstyle, caption={Lớp MetroFullTopo -- xây dựng topology modular}]
class MetroFullTopo(Topo):
    def build(self, **opts):
        self._build_backbone()   # CE/PE/P + P2P links
        self._build_branch1()    # Flat Network
        self._build_branch2()    # 3-Tier CDA
        self._build_branch3()    # Leaf-Spine
        self._build_glue_links() # CE <-> CE_LAN

    def _build_backbone(self):
        # CE routers (Customer Edge)
        self.addNode('ce1', cls=LinuxRouter, ip=None)
        self.addNode('ce2', cls=LinuxRouter, ip=None)
        self.addNode('ce3', cls=LinuxRouter, ip=None)
        # PE routers (Provider Edge)
        self.addNode('pe1', cls=LinuxRouter, ip=None)
        self.addNode('pe2', cls=LinuxRouter, ip=None)
        self.addNode('pe3', cls=LinuxRouter, ip=None)
        # P routers (Provider Core)
        self.addNode('p1', cls=LinuxRouter, ip=None)
        self.addNode('p2', cls=LinuxRouter, ip=None)
        self.addNode('p3', cls=LinuxRouter, ip=None)
        self.addNode('p4', cls=LinuxRouter, ip=None)
        # CE-PE links (point-to-point /30)
        self.addLink('ce1', 'pe1',
            intfName1='ce1-eth0', params1={'ip': '10.0.1.1/30'},
            intfName2='pe1-eth0', params2={'ip': '10.0.1.2/30'})
        # PE-P links
        self.addLink('pe1', 'p1',
            intfName1='pe1-eth1', params1={'ip': '10.10.11.1/30'},
            intfName2='p1-eth0',  params2={'ip': '10.10.11.2/30'})
        # P-P links (full-mesh P1/P2/P3/P4)
        self.addLink('p1', 'p2',
            intfName1='p1-eth1', params1={'ip': '10.20.12.1/30'},
            intfName2='p2-eth1', params2={'ip': '10.20.12.2/30'})
\end{lstlisting}



\subsection{Xây dựng chi nhánh 1 -- Flat Network}

\begin{lstlisting}[style=pythonstyle, caption={Hàm \_build\_branch1() -- Flat Network}]
def _build_branch1(self):
    """Switch access + 4 host + ce1_lan. Subnet: 10.1.0.0/24."""
    self.addNode('ce1_lan', cls=LinuxRouter, ip=None)
    s_acc1 = self.addSwitch('s_acc1',
                             cls=OVSSwitch,
                             failMode='standalone')
    # 4 hosts voi default route qua ce1_lan
    self.addHost('host1', ip='10.1.0.101/24',
                 defaultRoute='via 10.1.0.1')
    self.addHost('host2', ip='10.1.0.102/24',
                 defaultRoute='via 10.1.0.1')
    self.addHost('host3', ip='10.1.0.103/24',
                 defaultRoute='via 10.1.0.1')
    self.addHost('host4', ip='10.1.0.104/24',
                 defaultRoute='via 10.1.0.1')
    # host <-> switch access
    for h in ['host1','host2','host3','host4']:
        self.addLink(h, s_acc1)
    # switch access <-> ce1_lan (gateway LAN)
    self.addLink(s_acc1, 'ce1_lan',
        intfName2='ce1_lan-eth0',
        params2={'ip': '10.1.0.1/24'})
\end{lstlisting}

\subsection{Xây dựng chi nhánh 2 -- 3-Tier CDA}

\begin{lstlisting}[style=pythonstyle, caption={Hàm \_build\_branch2() -- 3-Tier Core-Distribution-Access}]
def _build_branch2(self):
    """3 lop: core1/core2 -> dist1/dist2 -> acc1/acc2/acc3 -> hosts."""
    self.addNode('ce2_lan', cls=LinuxRouter, ip=None)
    core1 = self.addSwitch('core1', cls=OVSSwitch, failMode='standalone')
    core2 = self.addSwitch('core2', cls=OVSSwitch, failMode='standalone')
    dist1 = self.addSwitch('dist1', cls=OVSSwitch, failMode='standalone')
    dist2 = self.addSwitch('dist2', cls=OVSSwitch, failMode='standalone')
    acc1 = self.addSwitch('acc1',  cls=OVSSwitch, failMode='standalone')
    acc2 = self.addSwitch('acc2',  cls=OVSSwitch, failMode='standalone')
    acc3 = self.addSwitch('acc3',  cls=OVSSwitch, failMode='standalone')
    self.addHost('admin1', ip='10.2.10.11/24',
                 defaultRoute='via 10.2.10.1')
    self.addHost('admin2', ip='10.2.10.12/24',
                 defaultRoute='via 10.2.10.1')
    self.addHost('lab1', ip='10.2.20.21/24',
                 defaultRoute='via 10.2.20.1')
    self.addHost('guest1', ip='10.2.30.31/24',
                 defaultRoute='via 10.2.30.1')
    # ce2_lan <-> Core (dual uplink)
    self.addLink('ce2_lan', core1,
        intfName1='ce2_lan-eth0', params1={'ip': '10.2.0.1/30'})
    self.addLink('ce2_lan', core2,
        intfName1='ce2_lan-eth1', params1={'ip': '10.2.0.5/30'})
    # Core <-> Distribution (cross-connected)
    self.addLink(core1, dist1); self.addLink(core1, dist2)
    self.addLink(core2, dist1); self.addLink(core2, dist2)
    # Distribution <-> Access
    self.addLink(dist1, acc1); self.addLink(dist1, acc2)
    self.addLink(dist2, acc2); self.addLink(dist2, acc3)
\end{lstlisting}

\subsection{Xây dựng chi nhánh 3 -- Leaf-Spine}

\begin{lstlisting}[style=pythonstyle, caption={Hàm \_build\_branch3() -- Leaf-Spine (2-tier)}]
def _build_branch3(self):
    """2 Spine + 3 Leaf (full-mesh) + 6 servers."""
    self.addNode('ce3_lan', cls=LinuxRouter, ip=None)
    spine1 = self.addSwitch('spine1', cls=OVSSwitch, failMode='standalone')
    spine2 = self.addSwitch('spine2', cls=OVSSwitch, failMode='standalone')
    leaf1  = self.addSwitch('leaf1',  cls=OVSSwitch, failMode='standalone')
    leaf2  = self.addSwitch('leaf2',  cls=OVSSwitch, failMode='standalone')
    leaf3  = self.addSwitch('leaf3',  cls=OVSSwitch, failMode='standalone')
    # Server zones
    self.addHost('web1', ip='10.3.10.11/24',
                 defaultRoute='via 10.3.10.1')
    self.addHost('dns1', ip='10.3.20.21/24',
                 defaultRoute='via 10.3.20.1')
    self.addHost('db1',  ip='10.3.30.31/24',
                 defaultRoute='via 10.3.30.1')
    # ce3_lan <-> Spine (dual uplink)
    self.addLink('ce3_lan', spine1,
        intfName1='ce3_lan-eth0', params1={'ip': '10.3.0.1/30'})
    self.addLink('ce3_lan', spine2,
        intfName1='ce3_lan-eth1', params1={'ip': '10.3.0.5/30'})
    # Spine <-> Leaf (full-mesh)
    for spine in [spine1, spine2]:
        for leaf in [leaf1, leaf2, leaf3]:
            self.addLink(spine, leaf)
    # Leaf <-> Servers
    self.addLink(leaf1, 'web1'); self.addLink(leaf1, 'web2')
    self.addLink(leaf2, 'dns1'); self.addLink(leaf2, 'dns2')
    self.addLink(leaf3, 'db1');  self.addLink(leaf3, 'db2')
\end{lstlisting}

\section{Cấu hình MPLS (Push -- Swap -- Pop)}

\subsection{Bước 1: Bật MPLS trên các node backbone}

\begin{lstlisting}[style=pythonstyle, caption={Hàm khởi tạo MPLS trên từng node}]
def configure_mpls(net):
    """Cau hinh Linux Kernel MPLS tren toan backbone."""
    MPLS_NODES = ['pe1', 'pe2', 'pe3', 'p1', 'p2', 'p3', 'p4']
    MPLS_INTERFACES = {
        'pe1': ['pe1-eth1', 'pe1-eth2'],
        'pe2': ['pe2-eth1', 'pe2-eth2'],
        'pe3': ['pe3-eth1', 'pe3-eth2'],
        'p1':  ['p1-eth0', 'p1-eth1', 'p1-eth2', 'p1-eth3'],
        'p2':  ['p2-eth0', 'p2-eth1', 'p2-eth2', 'p2-eth3'],
        'p3':  ['p3-eth0', 'p3-eth1', 'p3-eth2', 'p3-eth3', 'p3-eth4'],
        'p4':  ['p4-eth0', 'p4-eth1', 'p4-eth2', 'p4-eth3', 'p4-eth4'],
    }
    for name in MPLS_NODES:
        node = net[name]
        node.cmd('modprobe mpls_router   2>/dev/null || true')
        node.cmd('modprobe mpls_iptunnel 2>/dev/null || true')
        node.cmd('sysctl -w net.mpls.platform_labels=100000')
        for iface in MPLS_INTERFACES.get(name, []):
            node.cmd(f'sysctl -w net.mpls.conf.{iface}.input=1')
\end{lstlisting}

\subsection{Bước 2: Cấu hình PUSH tại PE routers (Ingress LER)}

\begin{lstlisting}[style=pythonstyle, caption={Cấu hình PUSH tại PE1, PE2, PE3}]
def _configure_pe_push(net):
    """PUSH: PE gan nhan MPLS vao goi IP den tu CE."""
    pe1, pe2, pe3 = net['pe1'], net['pe2'], net['pe3']

    # PE1: push label cho traffic den LAN2 (200) va LAN3 (300)
    pe1.cmd('ip route add 10.2.0.0/16  encap mpls 200 '
            'via 10.10.11.2 dev pe1-eth1')
    pe1.cmd('ip route add 10.3.0.0/16  encap mpls 300 '
            'via 10.10.11.2 dev pe1-eth1')

    # PE2: push label cho traffic den LAN1 (100) va LAN3 (310)
    pe2.cmd('ip route add 10.1.0.0/24  encap mpls 100 '
            'via 10.10.23.2 dev pe2-eth1')
    pe2.cmd('ip route add 10.3.0.0/16  encap mpls 310 '
            'via 10.10.24.2 dev pe2-eth2')

    # PE3: push label cho traffic den LAN1 (110) va LAN2 (210)
    pe3.cmd('ip route add 10.1.0.0/24  encap mpls 110 '
            'via 10.10.32.2 dev pe3-eth1')
    pe3.cmd('ip route add 10.2.0.0/16  encap mpls 210 '
            'via 10.10.32.2 dev pe3-eth1')
\end{lstlisting}

\subsection{Bước 3: Cấu hình SWAP tại P routers (Transit LSR)}

\begin{lstlisting}[style=pythonstyle, caption={Cấu hình SWAP LFIB tại P1, P2, P3, P4}]
def _configure_p_lfib(net):
    """SWAP: P nhan goi co nhan, thay bang nhan moi, day ra interface."""
    p1, p2, p3, p4 = net['p1'], net['p2'], net['p3'], net['p4']

    # ---- P1: SWAP ----
    p1.cmd('ip -f mpls route add 200 as 201 via inet 10.20.13.2 dev p1-eth2')
    p1.cmd('ip -f mpls route add 300 as 301 via inet 10.20.12.2 dev p1-eth1')
    p1.cmd('ip -f mpls route add 101 as 0   via inet 10.10.11.1 dev p1-eth0')
    p1.cmd('ip -f mpls route add 111 as 0   via inet 10.10.11.1 dev p1-eth0')

    # ---- P2: SWAP ----
    p2.cmd('ip -f mpls route add 110 as 111 via inet 10.20.12.1 dev p2-eth1')
    p2.cmd('ip -f mpls route add 210 as 211 via inet 10.20.24.2 dev p2-eth3')
    p2.cmd('ip -f mpls route add 301 as 0   via inet 10.10.32.1 dev p2-eth0')

    # ---- P3: SWAP ----
    p3.cmd('ip -f mpls route add 100 as 101 via inet 10.20.13.1 dev p3-eth2')
    p3.cmd('ip -f mpls route add 201 as 0   via inet 10.10.23.1 dev p3-eth1')
    p3.cmd('ip -f mpls route add 310 as 311 via inet 10.20.34.2 dev p3-eth4')

    # ---- P4: SWAP ----
    p4.cmd('ip -f mpls route add 311 as 0   via inet 10.10.34.1 dev p4-eth1')
    p4.cmd('ip -f mpls route add 211 as 0   via inet 10.10.24.1 dev p4-eth0')
\end{lstlisting}

\section{Hướng dẫn chạy topology}

\subsection{Khởi động topology tổng thể}

\begin{lstlisting}[style=bashstyle, caption={Chạy toàn bộ hệ thống Metro Ethernet MPLS}]
# Vao thu muc topology
cd ~/TKM2/topology

# Don dep Mininet cu (neu can)
sudo mn -c

# Chay topology tong the (3 chi nhanh + MPLS backbone)
sudo python3 metro_full.py
\end{lstlisting}

\subsection{Kiểm tra kết nối trong Mininet CLI}

\begin{lstlisting}[style=bashstyle, caption={Các lệnh kiểm tra cơ bản trong Mininet CLI}]
# Kiem tra tat ca node
mininet> nodes

# Ping tu chi nhanh 1 sang chi nhanh 2 (cross-branch)
mininet> host1 ping -c 5 10.2.10.11

# Ping tu chi nhanh 1 sang chi nhanh 3
mininet> host1 ping -c 5 10.3.10.11

# Xem bang dinh tuyen cua CE1
mininet> ce1 ip route

# Xem bang MPLS cua P1
mininet> p1 ip -f mpls route

# Traceroute tu host1 -> admin1 (xem duong qua backbone)
mininet> host1 traceroute -n 10.2.10.11
\end{lstlisting}

\subsection{Kết quả ping end-to-end thực tế -- Chi nhánh 1 (Flat)}

\begin{lstlisting}[style=termstyle, caption={Ping intra-branch: host1 $\to$ host2 và host1 $\to$ host3 (Flat Network)}]
>>> Test: host1->host2 [intra-flat]
  Run 1/3: ping 20 goi ICMP...  loss=0.0%  avg=0.066ms  min=0.030ms
  Run 2/3: ping 20 goi ICMP...  loss=0.0%  avg=0.046ms  min=0.030ms
  Run 3/3: ping 20 goi ICMP...  loss=0.0%  avg=0.045ms  min=0.030ms
  [OK] @ 10Mbps : tput=10.00Mbps  delay=0.0523ms  loss=0.0%   jitter=0.7588ms
  [OK] @ 50Mbps : tput=50.00Mbps  delay=0.0523ms  loss=0.0%   jitter=0.1565ms
  [OK] @ 100Mbps: tput=99.98Mbps  delay=0.0523ms  loss=0.0%   jitter=0.0686ms

>>> Test: host1->host3 [intra-flat]
  Run 1/3: ping 20 goi ICMP...  loss=0.0%  avg=0.065ms  min=0.035ms
  Run 2/3: ping 20 goi ICMP...  loss=0.0%  avg=0.046ms  min=0.035ms
  Run 3/3: ping 20 goi ICMP...  loss=0.0%  avg=0.047ms  min=0.035ms
  [OK] @ 10Mbps : tput=10.00Mbps  delay=0.0527ms  loss=0.0%   jitter=0.8368ms
  [OK] @ 50Mbps : tput=49.99Mbps  delay=0.0527ms  loss=0.0%   jitter=0.1626ms
  [OK] @ 100Mbps: tput=99.97Mbps  delay=0.0527ms  loss=0.0%   jitter=0.1868ms
\end{lstlisting}

\begin{lstlisting}[style=termstyle, caption={Ping cross-branch: host1 (Flat) $\to$ admin1 (3-Tier) và web1 (Leaf-Spine)}]
>>> Test: host1->admin1 [flat->3tier, cross-branch qua MPLS]
  Run 1/3: ping 20 goi ICMP...  loss=0.0%  avg=0.110ms  min=0.0797ms
  Run 2/3: ping 20 goi ICMP...  loss=0.0%  avg=0.100ms  min=0.0797ms
  Run 3/3: ping 20 goi ICMP...  loss=0.0%  avg=0.104ms  min=0.0797ms
  [OK] @ 10Mbps : tput=10.00Mbps  delay=0.1047ms  loss=0.0%    jitter=0.9715ms
  [OK] @ 50Mbps : tput=49.99Mbps  delay=0.1047ms  loss=0.0%    jitter=0.1969ms
  [OK] @ 100Mbps: tput=99.99Mbps  delay=0.1047ms  loss=0.0167% jitter=0.0893ms

>>> Test: host1->web1 [flat->leafspine, cross-branch qua MPLS]
  Run 1/3: ping 20 goi ICMP...  loss=0.0%  avg=0.122ms  min=0.0777ms
  Run 2/3: ping 20 goi ICMP...  loss=0.0%  avg=0.105ms  min=0.0777ms
  Run 3/3: ping 20 goi ICMP...  loss=0.0%  avg=0.103ms  min=0.0777ms
  [OK] @ 10Mbps : tput=10.00Mbps  delay=0.110ms  loss=0.0%  jitter=0.8527ms
  [OK] @ 50Mbps : tput=49.99Mbps  delay=0.110ms  loss=0.0%  jitter=0.1820ms
  [OK] @ 100Mbps: tput=99.99Mbps  delay=0.110ms  loss=0.0%  jitter=0.0725ms
\end{lstlisting}

\subsection{Kết quả ping end-to-end -- Chi nhánh 2 (3-Tier)}

\begin{lstlisting}[style=termstyle, caption={Ping intra-branch: admin1 $\to$ admin2 (same-VLAN), lab1, guest1 (inter-VLAN)}]
>>> Test: admin1->admin2 [same-VLAN10, intra-3tier]
  Run 1/3: ping...  loss=0.0%  avg=0.060ms  |  Run 2: avg=0.050ms  |  Run 3: avg=0.047ms
  [OK] @ 100Mbps: tput=99.98Mbps  delay=0.0523ms  loss=0.0%   jitter=0.1074ms

>>> Test: admin1->lab1 [inter-VLAN: VLAN10->VLAN20]
  Run 1/3: ping...  loss=0.0%  avg=0.081ms  |  Run 2: avg=0.075ms  |  Run 3: avg=0.071ms
  [OK] @ 10Mbps : tput=10.00Mbps  delay=0.0757ms  loss=0.0%  jitter=0.9006ms
  [OK] @ 50Mbps : tput=49.99Mbps  delay=0.0757ms  loss=0.0%  jitter=0.1311ms
  [OK] @ 100Mbps: tput=99.99Mbps  delay=0.0757ms  loss=0.0%  jitter=0.0331ms

>>> Test: admin1->guest1 [inter-VLAN: VLAN10->VLAN30]
  Run 1/3: ping...  loss=0.0%  avg=0.101ms  |  Run 2: avg=0.076ms  |  Run 3: avg=0.078ms
  [OK] @ 10Mbps : tput=10.00Mbps  delay=0.0850ms  loss=0.0%    jitter=0.8591ms
  [OK] @ 100Mbps: tput=99.99Mbps  delay=0.0850ms  loss=0.0067% jitter=0.0477ms
\end{lstlisting}

\begin{lstlisting}[style=termstyle, caption={Ping cross-branch: admin1 (3-Tier) $\to$ host1 (Flat) và web1 (Leaf-Spine)}]
>>> Test: admin1->host1 [3tier->flat, cross-branch qua MPLS]
  Run 1/3: ping...  avg=0.118ms  |  Run 2: avg=0.100ms  |  Run 3: avg=0.104ms
  [OK] @ 100Mbps: tput=99.98Mbps  delay=0.1073ms  loss=0.0%  jitter=0.1518ms
  TCP throughput avg: 24333.571 Mbps

>>> Test: admin1->web1 [3tier->leafspine, cross-branch qua MPLS]
  Run 1/3: ping...  avg=0.126ms  |  Run 2: avg=0.106ms  |  Run 3: avg=0.108ms
  [OK] @ 100Mbps: tput=99.99Mbps  delay=0.1133ms  loss=0.0%  jitter=0.1502ms
  TCP throughput avg: 22127.227 Mbps
\end{lstlisting}

\subsection{Kết quả ping end-to-end -- Chi nhánh 3 (Leaf-Spine)}

\begin{lstlisting}[style=termstyle, caption={Ping intra-branch: web1 $\to$ web2 (same-leaf), dns1, db1 (2-hop)}]
>>> Test: web1->web2 [same-leaf1, intra-leafspine]
  Run 1/3: ping...  avg=0.067ms  |  Run 2: avg=0.047ms  |  Run 3: avg=0.047ms
  [OK] @ 100Mbps: tput=99.99Mbps  delay=0.0537ms  loss=0.0%  jitter=0.0397ms

>>> Test: web1->dns1 [leaf1->leaf2, 2-hop via Spine]
  Run 1/3: ping...  avg=0.076ms  |  Run 2: avg=0.070ms  |  Run 3: avg=0.068ms
  [OK] @ 10Mbps : tput=10.00Mbps  delay=0.0713ms  loss=0.0%   jitter=0.7190ms
  [OK] @ 100Mbps: tput=99.98Mbps  delay=0.0713ms  loss=0.07%  jitter=0.1095ms

>>> Test: web1->db1 [leaf1->leaf3, 2-hop via Spine]
  Run 1/3: ping...  avg=0.087ms  |  Run 2: avg=0.072ms  |  Run 3: avg=0.073ms
  [OK] @ 100Mbps: tput=99.99Mbps  delay=0.0767ms  loss=0.0%  jitter=0.0813ms
\end{lstlisting}

\begin{lstlisting}[style=termstyle, caption={Ping cross-branch: web1 (Leaf-Spine) $\to$ host1 (Flat) và admin1 (3-Tier)}]
>>> Test: web1->host1 [leafspine->flat, cross-branch qua MPLS]
  Run 1/3: ping...  avg=0.118ms  |  Run 2: avg=0.106ms  |  Run 3: avg=0.107ms
  [OK] @ 100Mbps: tput=99.99Mbps  delay=0.1103ms  loss=0.0267% jitter=0.1021ms
  TCP throughput avg: 22105.305 Mbps

>>> Test: web1->admin1 [leafspine->3tier, cross-branch qua MPLS]
  Run 1/3: ping...  avg=0.125ms  |  Run 2: avg=0.111ms  |  Run 3: avg=0.107ms
  [OK] @ 100Mbps: tput=99.99Mbps  delay=0.1143ms  loss=0.0967% jitter=0.0915ms
  TCP throughput avg: 22208.975 Mbps

=== TONG KET: 42 records, 14 kich ban, lap lai 3 lan ===
  Tat ca cross-branch: loss < 0.1%, delay < 0.12ms, tput > 99.97%
  CSV: results/csv/results_20260418_010018.csv
\end{lstlisting}

\subsection{Đo thông lượng iperf3 trong CLI}

\begin{lstlisting}[style=bashstyle, caption={Đo thông lượng TCP và UDP bằng iperf3}]
# Do TCP throughput: host1 -> admin1 (cross-branch)
mininet> admin1 iperf3 -s &
mininet> host1  iperf3 -c 10.2.10.11 -t 10

# Do UDP jitter va packet loss: host1 -> web1
mininet> web1   iperf3 -s &
mininet> host1  iperf3 -c 10.3.10.11 -u -b 100M -t 10

# Thoat va don dep
mininet> exit
sudo mn -c
\end{lstlisting}


\section{Bộ công cụ đo lường tự động}

\subsection{Cấu hình kịch bản đo lường}

\begin{lstlisting}[style=pythonstyle, caption={Cấu hình kịch bản đo lường trong tests.yaml}]
scenarios:
  - name: flat
    arch_name: "Flat Network (Chi nhanh 1)"
    tests:
      - src: host1
        dst_ip: 10.1.0.102
        description: "host1->host2 [intra-flat]"
        type: intra
      - src: host1
        dst_ip: 10.2.10.11
        description: "host1->admin1 [flat->3tier]"
        type: cross

measurement:
  ping_count: 20          # So goi ICMP
  iperf_duration: 10      # Thoi gian do (giay)
  bandwidths: [10M, 50M, 100M]  # Cac muc tai UDP
  repeats: 3              # So lan lap lai moi kich ban
\end{lstlisting}

\begin{lstlisting}[style=bashstyle, caption={Chạy bộ đo lường tự động}]
# Chay do luong va luu ket qua
cd ~/TKM2
sudo python3 scripts/run_measurements.py \
    --topology topology/metro_full.py \
    --config   config/tests.yaml \
    --output   results/

# Ket qua luu tai:
#   results/csv/results_<timestamp>.csv
#   results/raw/all_<timestamp>/

# Ve bieu do tu CSV
python3 scripts/plot_results.py \
    --csv results/csv/results_20260418_010018.csv
\end{lstlisting}
